\documentclass[11pt]{article}
\usepackage{amsmath,graphicx}
\title{Entropy-Guided Adaptive Token-Budget Routing for Long-Context Transformers}

\begin{document}
\maketitle

\begin{abstract}
Understanding how attention mass concentrates across ultra-long sequences is of paramount importance for scaling modern Transformer inference. However, existing KV-Cache eviction paradigms exhibit fundamental limitations because fixed-window pruning fails to retain transient long-range retrieval anchors. In this study, we delve into an entropy-guided adaptive routing (AdaptRoute) framework that leverages layer-wise attention priors to guide dynamic block selection. Across a benchmark of 88 reasoning tasks ($n = 88$, 128K context), our framework reduces perplexity gap (PPL) from 0.412 to 0.286 nats ($p < 0.001$), offering a comprehensive and auditable blueprint for memory-constrained LLM serving.
\end{abstract}

\section{Introduction}
\label{sec:intro}

Transformer attention weights span orders of magnitude across early syntactic layers and deep retrieval heads~\cite{vaswani2017neurips,dao2022neurips}. Moreover, tight GPU memory budgets make exhaustive KV-Cache retention infeasible in real-world deployments. Traditional eviction protocols rely on uniform or sliding token windows, meaning a large fraction of memory budget is spent on low-entropy filler tokens while missing long-range reasoning anchors, as illustrated in Fig.~\ref{fig:profile_overview}.

Furthermore, purely heuristic pruning methods struggle under tight cache budgets ($N \le 8$ blocks per head) because they lack prior knowledge of head specialization and chain-of-thought dependency spans~\cite{wei2022neurips}. Notably, coupling a reduced-order entropy prior with sequential expected information gain (EIG) selection may suggest a practical path toward higher accuracy without increasing GPU memory.

\section{Methodology}
\label{sec:method}

\subsection{Hierarchical Attention Entropy Prior}
We model block importance $c(z) = \log_{10} C(z)$ over block index $z \in [0, 500]$ as a Gaussian process whose mean function $m_\phi(z)$ is anchored by a layer-wise attention balance:

\begin{equation}
\label{eq:transport}
\frac{\partial}{\partial z}\!\left(w_s(z)\, C(z) - K_z(z)\frac{\partial C(z)}{\partial z}\right) = S_{\text{head}}(z) - \lambda_{\text{decay}} C(z),
\end{equation}

where $w_s(z)$ denotes the positional decay velocity, $K_z(z)$ is cross-block diffusion, and the prior variance satisfies $\sigma_z^2 = 0.64$.

\subsection{Sequential Expected Information Gain}
At each decode step $t$, the AdaptRoute policy selects the next key--value block $z_t^\star$ by maximizing the mutual information between the block logit $y_t$ and the full attention output $\mathbf{c}$:

\begin{equation}
\label{eq:eig}
z_t^\star = \arg\max_{z \in \mathcal{Z}} \; \mathbb{E}_{p(y \mid z, \mathcal{D}_{t-1})}\!\left[ D_{\text{KL}}\!\left(p(\mathbf{c} \mid \mathcal{D}_{t-1}, z, y) \,\|\, p(\mathbf{c} \mid \mathcal{D}_{t-1})\right) \right].
\end{equation}

This paragraph describes our operational GPU memory cost model: each deep layer above 200 blocks incurs a bandwidth penalty proportional to cache depth, $\eta_{\text{cost}}(z) = 1 + 0.36 (z / 500)$, ensuring that distant blocks are queried only when expected distortion reduction outweighs latency cost.

\section{Experiments and Results}
\label{sec:results}

We evaluate AdaptRoute against uniform pruning, sliding-window eviction, and unconstrained top-k caching across 88 benchmark tasks summarized in Table~\ref{tab:ablation}. Under an 8-block budget, AdaptRoute achieves a PPL gap of 0.286 compared to 0.412 for uniform eviction (a 30.6\% reduction, $p < 0.001$) and improves long-range needle retrieval accuracy from 61.5\% to 84.2\%.

\begin{figure}[t]
\caption{Comparison of long-context reasoning accuracy under an 8-block KV-Cache budget across uniform, sliding-window, and entropy-guided AdaptRoute policies.}
\label{fig:profile_overview}
\end{figure}

\begin{table}[t]
\caption{Quantitative ablation of entropy prior, cost weighting, and block budget across $n = 88$ tasks.}
\label{tab:ablation}
\begin{tabular}{lccc}
\toprule
Strategy & Block Budget & PPL Gap & Retrieval Acc. (\%) \\
\midrule
Uniform & 8 & 0.412 & 61.5 \\
Sliding Window & 8 & 0.368 & 68.4 \\
AdaptRoute (Ours) & 8 & \textbf{0.286} & \textbf{84.2} \\
\bottomrule
\end{tabular}
\end{table}

\section{Discussion}
\label{sec:discussion}

Our results indicate that embedding layer-wise entropy structure into sequential block routing substantially narrows output distortion at 128K sequence lengths. Under the stated single-host serving assumptions, the method provides calibrated latency budgets suitable for practical LLM deployment.

\end{document}
